
\section{Limitations}
\label{sec:lim}

\paragraph{This is predictive performance, not value}
Nothing here measures what a register costs to acquire or maintain, whether
acting on a prediction changes an outcome, what a false review costs the
organisation whose data these are, or what a data programme is worth. The
quantity is an increment in predictive performance, read at a declared
exchange rate in Section~\ref{sec:dca}. A reader who converts
Table~\ref{tab:master} into a procurement decision is doing something this
evidence does not support.

\paragraph{The design space is declared, not complete}
$\mathcal{S}$ has \nAxes\ axes here. It could have more: hyperparameter
tuning protocol, missing-data treatment for the baseline fields, the horizon
of the target, the label-noise model. The cohort and eligibility rule is
varied only on the case study's log, where the two cohorts are reported side
by side; the rest are fixed. A surface is only ever a surface over what its
author declared, and the honest form of the claim is ``the increment moves
this much over \emph{this} space'', not ``over all defensible choices''.

\paragraph{The grid is not uniform across the corpus, and the bootstrap's is
smaller than the surface's} The learner axis carries four levels on the case
study's log and two elsewhere, and the nested bootstrap covers four
register-quality conditions where the surface covers \nQuality. Both are
compute budgets, both were declared before the runs rather than chosen after
them, and both mean that a sensitivity index computed on one pair is not
strictly comparable with one computed on another. The alternative --- a
uniform grid at the size the smallest budget allows --- would have cost more
information than it bought.

\paragraph{Two of the new objects have stated scope conditions} The
decomposition inherits the metric's scale (Proposition~\ref{prop:scale}), so
``which choice matters most'' is a statement about a scale as well as about a
surface. And the regret benchmark assumes the reader's own specification is
drawn from the declared admissible set under a stated weighting
(Proposition~\ref{prop:optimal}); two weightings are reported, and a reader
whose practice matches neither will get a different number.

\paragraph{One estate for the case study}
The configuration management database analysis is one bank's estate, on one
public log, and the availability question at its centre is settled by
probability rather than by a timestamped field-value history. The strongest
version of this study would have an organisational partner supplying
timestamped field histories, a known prediction time, historical register
population and accuracy, operational review costs, and practitioner
confirmation of which fields are admissible. We do not have one, and
Section~\ref{sec:cmdb} is framed as a benchmark case in consequence.

\paragraph{The corpus is heterogeneous, and here is how much}
The \nPairs\ pairs are not replications of one another. The registers, the
free fields and the targets mean different things across domains, and the
claim deserves a number rather than a hedge: across the \nLogs\ logs the two
registered targets label the same case identically on a median
\targetAgreement\ of cases, ranging from \targetAgreementLo\ to
\targetAgreementHi. They are two questions, not one question asked twice ---
which is what makes the target a real axis, and what makes averaging across
pairs meaningless. The corpus's handover target and the case study's
reassignment target agree on \crossTargetAgreement\ of the incidents the two
cohorts share. Section~\ref{sec:multilog} therefore reports a benchmark of
specification sensitivity across heterogeneous problems, and does not average
across them.

\paragraph{The quality mechanisms are synthetic}
Population loss, identity error, reconciliation failure and staleness are
constructed from the training half, not observed. A real register degrades in
ways correlated with the outcome --- the items nobody maintains are the items
nobody escalates --- and a mechanism built to be outcome-free by construction
cannot capture that. Validating the mechanisms against a genuinely incomplete
or historically evolving register is the obvious next study and we do not have
the data for it.

\paragraph{Inference still rests on assumptions}
The nested bootstrap refits the pipeline and the blocks carry temporal
dependence, but the split point is fixed within a draw, the block length is a
rule of thumb, and the max-$t$ band's family-wise coverage is asymptotic. The
simulation of Section~\ref{sec:sim} measures what this costs in \nWorlds\ worlds;
it does not certify the intervals on a seventh.

\paragraph{The practice pilot is a pilot}
Section~\ref{sec:practice} states its own sampling error, coder error and
resolution, and none of the paper's claims rests on it.

\section{Conclusion}
\label{sec:conc}

The incremental predictive performance of a recorded field is a functional of
design choices, not a number. This paper made that operational: a declared
design space, a surface over it, a decomposition of the surface's variance
into the contribution of each choice, simultaneous confidence bands from a
bootstrap that refits the model, region labels that say when a sign is robust,
and a decision-theoretic benchmark that measures what one number gets wrong.

Four results are worth carrying away.

\textbf{The choices move the answer, and no one of them is the culprit.}
Varying the baseline alone moves the increment by \spreadMedian\ AUC at the
median log--target pair. But no single axis dominates the disagreement: the
largest median first-order index is \sobolMaxPct, the interactions carry
\sobolInteractionPct, and which axis leads is a property of the pair. An
analyst cannot learn from somebody else's paper which of their own choices
their answer turns on. They have to compute it.

\textbf{The sign is a property of the specification, not of the data.} On
\nSignVaries\ of \nPairs\ pairs, between a tenth and nine tenths of the
admissible cells are positive, and a conventional one-number report misstates
the sign for a median \misreportMedianPct\ of the rest.

\textbf{There is a defensible collapse, and it is not the one anybody uses.}
Proposition~\ref{prop:optimal} identifies it --- the weighted mean increment
over the admissible set --- and on this corpus the conventional report costs
\regretOneNumberExcessPct\ more expected regret than it does, while the
cautious rule of adopting only where the whole surface agrees costs more than
either. The collapse depends on which specifications the reader's own practice
makes likely, which the analyst does not know. So: collapse for a decision,
report the surface.

\textbf{Availability is an argument of the estimand, not a preliminary.} On
the case study the register is worth \VTone\ at one decision time and
\VTtwoKnow\ at another, and no evidence we have settles which one a reader
should stand at. We report both, and a curve between them, rather than
choosing --- because choosing is what produced the earlier version's headline
and then removed it.

The practical recommendation is narrow and cheap once the surface exists:
report the surface, the decomposition, the region and $\rho$; make the
absolute increments primary; use a ratio last, with a Fieller set, and only
where its denominator is resolvably positive. \texttt{fieldvalue} does all of
it from a data frame, and refuses, by design, to give back a single number.
